\tcbset{
  scriptbox/.style={
    enhanced,
    breakable=false,
    colback=LectureBoxBackground,
    colbacktitle=LectureBoxBackground,
    colframe=LectureBoxBackground,
    boxrule=0pt,
    borderline west={1.2pt}{0pt}{THMBlue},
    sharp corners,
    left=2.5mm,right=2mm,top=1.5mm,bottom=1.5mm,
    before skip=1.0\baselineskip,
    after skip=1.0\baselineskip,
    fonttitle=\bfseries,
    coltitle=black
  }
}

%---------
%--Boxes--
%---------

% Definition
\newtcbtheorem[number within=chapter]{DefinitionBox}{\LectureDefinitionLabel}
  {scriptbox,breakable=false}{def}

% Function properties are numbered chronologically across the script.
\newtcbtheorem{FunctionPropertyBox}{\Lang{Funktionseigenschaft}{Function property}}
  {scriptbox,breakable=false}{def}

% Merke
\newtcbtheorem[number within=chapter]{KeyPointBox}{\Lang{Merke}{Key Point}}
  {scriptbox,borderline west={2pt}{0pt}{THMBlue}}{key}

% Beispiel
\newtcbtheorem[number within=chapter]{ExampleBox}{\LectureExampleLabel}
  {scriptbox,breakable=false,borderline west={1.2pt}{0pt}{THMBlack}}{ex}

% Theorem
\newtcbtheorem[number within=chapter]{TheoremBox}{\LectureTheoremLabel}
  {scriptbox,borderline west={1.8pt}{0pt}{THMBlue}}{thm}

% Remarks have their own chapter counter, including individually titled remarks.
\tcbset{
  remark title/.style={title={\Lang{Hinweis}{Remark}~\thetcbcounter: #1}}
}
\newtcolorbox[auto counter,number within=chapter]{RemarkBox}[1][]{
  scriptbox,
  borderline west={1.2pt}{0pt}{THMBlack},
  title={\Lang{Hinweis}{Remark}~\thetcbcounter},#1
}
